\clearpage
\hypertarget{native-270-pure17}{}\label{native-270-pure17}
\Article{17}
For a clearer illustration, let us return to the example of Article~13, where we found
\[
\begin{aligned}
A'=A'''&=2+4R+R^3+2R^4=2R-2R^2-R^3\\
A''&=2+R+4R^2+2R^3=-R+2R^2-2R^4
\end{aligned}
\]
Taking $R$ to have the value $\cos72^0+i\sin72^0$, we obtain
\[
\begin{aligned}
A'=A'''&=2\cos72^0-3\cos144^0+i(2\sin72^0-\sin144^0)\\
A''&=-3\cos72^0+2\cos144^0+i(\sin72^0+2\sin144^0)
\end{aligned}
\]
The angles $b',b''$ will therefore be determined by the equations
\[
\begin{aligned}
1)\qquad\sin b'&=\frac{2\sin72^0-\sin144^0}{\surd11}\\
2)\qquad\cos b'&=\frac{2\cos72^0-3\cos144^0}{\surd11}\\
3)\qquad\operatorname{tang}b'&=\frac{2\sin72^0-\sin144^0}{2\cos72^0-3\cos144^0}\\
4)\qquad\sin b''&=\frac{\sin72^0+2\sin144^0}{\surd11}\\
5)\qquad\cos b''&=\frac{-3\cos72^0+2\cos144^0}{\surd11}\\
6)\qquad\operatorname{tang}b''&=\frac{\sin72^0+2\sin144^0}{-3\cos72^0+2\cos144^0}
\end{aligned}
\]
Any one of the equations $1,2,3$ suffices to determine the angle $b'$, provided the quadrant in which it is to be taken is known; this must be decided from the signs of the quantities $2\sin72^0-\sin144^0$, $2\cos72^0-3\cos144^0$. The same holds for the angle $b''$. In our case $b'$ is to be taken between $0$ and $90^0$, and $b''$ between $90^0$ and $180^0$. If the numerator and denominator of equation $3$ are multiplied by $-3\cos72^0+2\cos144^0$, it becomes
\SourceLeaf{271}
\[
\operatorname{tang}b'=\frac2{31}\{-\sin72^0+13\sin144^0\}
\]
Similarly, from equation $6$, on multiplying the numerator and denominator by $2\cos72^0-3\cos144^0$, we obtain
\[
\operatorname{tang}b''=\frac2{31}\{-13\sin72^0-\sin144^0\}
\]
Numerically, this gives
\[
\begin{aligned}
\operatorname{tang}b'&=+0{,}4316226944,\quad\log\operatorname{tang}b'=9{,}6351042715\\
&\hspace{1em}b'=23^020'46''04603\\
\operatorname{tang}b''&=-0{,}8355819332,\quad\log\operatorname{tang}b''=9{,}9219890411\mathrm n\\
&\hspace{1em}b''=140^07'6''52441
\end{aligned}
\]
from which we derive
\[
5f'=186^048'38''61647,\qquad f'=37^021'43''723294
\]

.We therefore have
\[
\begin{aligned}
(2,1)&=-\frac15+\frac{\surd11}5\{2\cos37^021'43''723294+2\cos51^022'41''400558\}\\
(2,2)&=-\frac15+\frac{\surd11}5\{2\cos325^021'43''723294+2\cos267^022'41''400558\}\\
(2,4)&=-\frac15+\frac{\surd11}5\{2\cos253^021'43''723294+2\cos123^022'41''400558\}\\
(2,8)&=-\frac15+\frac{\surd11}5\{2\cos181^021'43''723294+2\cos339^022'41''400558\}\\
(2,5)&=-\frac15+\frac{\surd11}5\{2\cos109^021'43''723294+2\cos195^022'41''400558\}
\end{aligned}
\]
from which we find
\[
\begin{array}{rclcl}
(2,1)&=&+1{,}6825070652&=&2\cos\frac{360^0}{11}\\[3pt]
(2,2)&=&+0{,}8308299&=&2\cos\frac{720^0}{11}\\[3pt]
(2,4)&=&&=&2\cos\frac{1440^0}{11}\\[3pt]
(2,8)&=&&=&2\cos\frac{2880^0}{11}\\[3pt]
(2,5)&=&&=&2\cos\frac{1800^0}{11}
\end{array}
\]

\Article{18}
The equation $x^{17}-1=0$, which we already treated by a different method in the \emph{Disquis. Arithm.}, will provide another example. We shall therefore set $n=17$, $\beta=8$, $g=3$; the following indices then correspond to the numbers:
\SourceLeaf{272}
\[
\begin{array}{r*{16}{@{\;}r}}
\text{numbers}&1\pdot&2\pdot&3\pdot&4\pdot&5\pdot&6\pdot&7\pdot&8\pdot&9\pdot&10\pdot&11\pdot&12\pdot&13\pdot&14\pdot&15\pdot&16\\
\text{indices}&0\pdot&14\pdot&1\pdot&12\pdot&5\pdot&15\pdot&11\pdot&10\pdot&2\pdot&3\pdot&7\pdot&13\pdot&4\pdot&9\pdot&6\pdot&8
\end{array}
\]
From this we find
\[
\begin{aligned}
A'=A''''''&=2R+2R^2+3R^4+4R^5+2R^6+2R^7\\
A''=A'''''&=2+3R+R^3+R^4+3R^5+4R^6+R^7\\
A'''=A''''&=3+3R+2R^2+3R^3+R^5+2R^6+R^7
\end{aligned}
\]
or, since in this case $R^4+1=0$,
\[
\begin{aligned}
A'=A''''''&=-3-2R-2R^3\\
A''=A'''''&=1-4R^2\\
A'''=A''''&=3+2R+2R^3
\end{aligned}
\]
Putting $R=\cos45^0+i\sin45^0$, we therefore have
\[
\begin{aligned}
A'=A''''''&=-3-2i\surd2,\\
A''=A'''''&=1-4i,\\
A'''=A''''&=3+2i\surd2
\end{aligned}
\]

Thus $b',b'',b'''$ will be found from the equations
\[
\begin{aligned}
\sin b'&=-\surd\tfrac8{17}&\sin b''&=-\surd\tfrac{16}{17}&\sin b'''&=+\surd\tfrac8{17}\\
\cos b'&=-\surd\tfrac9{17}&\cos b''&=+\surd\tfrac1{17}&\cos b'''&=+\surd\tfrac9{17}\\
\operatorname{tang}b'&=+\surd\tfrac89&\operatorname{tang}b''&=-4&\operatorname{tang}b'''&=+\surd\tfrac89
\end{aligned}
\]
from which we deduce
\[
b'=223^018'49'',\quad b''=284^02'10,\quad b'''=43^018'49''=b'-180^0
\]
\[
4f'=550^039'48'',\qquad f'=137^039'57''
\]
\[
\!\begin{aligned}
(2,1)&=\mkern-2mu-\frac18+\frac{\surd17}8\{2\cos137^039'57''+2\cos52^01'5''+2\cos265^038'52''+1\}\\
(2,4)&=\mkern-2mu-\frac18+\frac{\surd17}8\{2\cos92^039'57''+2\cos322^01'5''+2\cos130^038'52''-1\}\\
(2,9)&=\mkern-2mu-\frac18+\frac{\surd17}8\{2\cos47^039'57''+2\cos232^01'5''+2\cos355^038'52''+1\}\\
(2,10)&=\mkern-2mu-\frac18+\frac{\surd17}8\{2\cos2^039'57''+2\cos142^01'5''+2\cos220^038'52''-1\}\\
(2,13)&=\mkern-2mu-\frac18+\frac{\surd17}8\{2\cos317^039'57''+2\cos52^01'5''+2\cos85^038'52''+1\}\\
(2,5)&=\mkern-2mu-\frac18+\frac{\surd17}8\{2\cos272^039'57''+2\cos322^01'5''+2\cos310^038'52''-1\}\\
(2,15)&=\mkern-2mu-\frac18+\frac{\surd17}8\{2\cos227^039'57''+2\cos232^01'5''+2\cos175^038'52''+1\}\\
(2,11)&=\mkern-2mu-\frac18+\frac{\surd17}8\{2\cos182^039'57''+2\cos142^01'5''+2\cos40^038'52''-1\}
\end{aligned}
\]

\SourceLeaf{273}
\[
\begin{array}{rclcl}
\frac12(2,1)&=&+0{,}092268&=&\cos\frac4{17}360^0\\[3pt]
\frac12(2,3)&=&&=&\cos\frac5{17}360^0\\[3pt]
\frac12(2,9)&=&&=&\cos\frac2{17}360^0\\[3pt]
\frac12(2,10)&=&&=&\cos\frac6{17}360^0\\[3pt]
\frac12(2,13)&=&+0{,}93247&=&\cos\frac1{17}360^0\\[3pt]
\frac12(2,5)&=&&=&\cos\frac3{17}360^0\\[3pt]
\frac12(2,15)&=&&=&\cos\frac8{17}360^0\\[3pt]
\frac12(2,11)&=&&=&\cos\frac7{17}360^0
\end{array}
\]
\begin{center}*\qquad\qquad*\\*\end{center}

From these more general investigations of the functions $[r,R]$, which further clarify and extend the second theory of pure equations begun in Article~360 of the \emph{Disquiss. Ar.}, we proceed to consider certain special cases more closely, namely those in which definite values are taken for $\beta$. Several investigations, as fruitful as they are elegant, will emerge from this; some had already been treated in the \emph{Disq. Ar.} (Articles\ .\ .\ .), though by a different method, while others are to be regarded as entirely new. We reserve for another paper, to be published soon, the development of the remarkable connection between these investigations and higher arithmetic, which thereby makes very great and hitherto unexpected advances. — Moreover, throughout the following investigation we shall assume that $r$ is taken to be a \emph{proper} root of the equation $x^n-1=0$, and $R$ a proper root of the equation $R^\beta-1=0$.

\Article{19}
We begin with the value $\beta=2$, so that $R$ must take the value $-1$. Our function $[r,R]$ therefore becomes
\[
=r-r^g+r^{g^2}-r^{g^3}\dotsfour-r^{g^{n-2}}
\]
and we have
\[
[r,R]=-[r^g,R]=+[r^{g^2},R]=-[r^{g^3},R]\quad\text{etc.}
\]
and, in general, if $\lambda$ denotes any integer not divisible by $n$,
\SourceLeaf{274}
\[
\begin{aligned}
\relax[r^\lambda,R]&=+[r,R]\quad\text{if }\lambda\text{ is a quadratic residue of }n,\\
\relax[r^\lambda,R]&=-[r,R]\quad\text{if }\lambda\text{ is a quadratic nonresidue of }n.
\end{aligned}
\]
\Needspace{11\baselineskip}
Further, if the quadratic residues of $n$ among $1,2,3\dots n-1$ are denoted generically by $a$, and the nonresidues of $n$ within the same limits by $b$, it is evident that the numbers
\[
1,g^2,g^4\dotsfour g^{n-3}
\]
with their order disregarded, are congruent modulo $n$ to the numbers $a$, and similarly the numbers
\[
g,g^3,g^5\dotsfour g^{n-2}
\]
are congruent to $b$, so that $[r,R]=\Sigma r^a-\Sigma r^b$.

If, therefore, we put $\frac{360^0}n=\omega$ and $r=\cos k\omega+i\sin k\omega$, then $[r,R]=\Sigma\cos ak\omega-\Sigma\cos bk\omega+i\Sigma\sin ak\omega-i\Sigma\sin bk\omega$. Now, by Article~14, the square of the function $[r,R]$ will be $=+n$ or $=-n$ according as $n$ has the form $4z+1$ or $4z-1$, and hence in the former case $[r,R]=\pm\surd n$, in the latter $[r.R]=\pm i\surd n$; the sign placed before the radical quantity, however, remains ambiguous. From this we obtain the following sums:

\noindent I.\quad If $n$ has the form $4z+1$,
\[
\begin{aligned}
\Sigma\cos ak\omega-\Sigma\cos bk\omega&=\pm\surd n\\
\Sigma\sin ak\omega-\Sigma\sin bk\omega&=0
\end{aligned}
\]
\noindent II.\quad If $n$ has the form $4z-1$,
\[
\begin{aligned}
\Sigma\cos ak\omega-\Sigma\cos bk\omega&=0\\
\Sigma\sin ak\omega-\Sigma\sin bk\omega&=\pm\surd n
\end{aligned}
\]

Moreover, since the whole collection of numbers $a,b$ evidently coincides with $1,2,3\dots n-1$, we have $\Sigma r^a+\Sigma r^b=r+r^2+r^3+\text{etc.}+r^{n-1}=-1$, and therefore $\Sigma\cos ak\omega+\Sigma\cos bk\omega=-1$, $\Sigma\sin ak\omega+\Sigma\sin bk\omega=0$. These identities and the preceding sums give:

\noindent I.\quad For the former case,
\[
\begin{aligned}
\Sigma\cos ak\omega&=-\tfrac12\pm\tfrac12\surd n\\
\Sigma\cos bk\omega&=-\tfrac12\mp\tfrac12\surd n\\
\Sigma\sin ak\omega&=\Sigma\sin bk\omega=0
\end{aligned}
\]

\SourceLeaf{275}
\noindent II.\quad For the latter case,
\[
\begin{aligned}
\Sigma\cos ak\omega&=\Sigma\cos bk\omega=-\tfrac12\\
\Sigma\sin ak\omega&=\pm\tfrac12\surd n\\
\Sigma\sin bk\omega&=\mp\tfrac12\surd n
\end{aligned}
\]
\Needspace{9\baselineskip}
These sums were already obtained by a method not very different in Article~356 of the \emph{Disquiss. Arr.}; neither method can remove the ambiguity in the sign to be placed before the radical quantity. We have, however, recently overcome this shortcoming in a separate paper, where it was proved that, for the value $k=1$, the upper signs must be taken in all the formulas given.
\BookRule{.3\textwidth}
